\begin{examplebox}
\textbf{\textcolor{CExample}{例 1（初中起点：正比例函数的对称轴）}}

\begin{center}
\begin{tikzpicture}[font=\small,>=Stealth,line cap=round,line join=round]
\path[use as bounding box] (-6.5,-3.55) rectangle (6.5,3.55);
\draw[black!20,rounded corners=3pt] (-6.5,-3.55) rectangle (6.5,3.55);
\node[font=\small\bfseries] at (0,3.08) {完整直线：自身与每条垂线都是对称轴};

\begin{scope}[rotate=63.43495]
\draw[CExplanation,very thick] (-2.6,0)--(2.6,0);
\foreach \u in {-1.2,0,1.2} {\draw[CProof,dashed] (\u,-1.5)--(\u,1.5);}
\draw[CStatement,<->,thick] (-.9,0)--(.9,0);
\fill[CStatement] (-.9,0) circle(2pt); \fill[CStatement] (.9,0) circle(2pt);
\draw (0,.23)--(.23,.23)--(.23,0);
\end{scope}
\node[anchor=west] at (2.2,2.1) {蓝色实线：$y=2x$};
\node[anchor=west] at (2.2,1.35) {虚线：三条对称轴};
\node[anchor=west] at (2.2,.6) {这样的轴有无数条};

\node[align=center] at (0,-3.05) {关于自身，每个点不动；关于任一垂线，直线上的点成对交换。};
\end{tikzpicture}
\end{center}


\textbf{题目：} $y=2x$ 只有一条对称轴吗？判断 $y=2x$、$y=-\frac12x$、$x=0$ 是否是其对称轴。

\textbf{解：} 曲线上每点都在 $y=2x$ 这条直线上，关于自身对称后保持不动，故它是对称轴。对于 $y=-\frac12x$，它经过原点并且与 $y=2x$ 垂直；在原直线上取 $(t,2t)$，其关于原点另一侧的 $(-t,-2t)$ 与它关于垂线 $y=-\frac12x$ 对称（原点是中点且连线垂直于该轴），故它也是对称轴。不仅如此，\textbf{每条过 $y=2x$ 上一点且垂直于它的直线}都是其对称轴。反之，$(1,2)$ 关于 $x=0$ 的对称点 $(-1,2)$ 不满足 $y=2x$，因此 $x=0$ 不是对称轴。

\textbf{结论：} \textbf{一条完整直线具有无数条对称轴}；不要将一般曲线的“唯一轴”直觉套到直线上。

\medskip
\textbf{\textcolor{CExample}{例 2（反比例函数：两条渐近线与两条对称轴）}}

\textbf{题目：} 证明 $y=\dfrac1x$ 关于 $y=x$、$y=-x$ 对称，并判断 $x=0$ 是否对称轴。

\textbf{解：} 化成 $xy=1$（$x\ne0$）。代换 $(x,y)\mapsto(y,x)$ 得 $yx=1$；代换 $(x,y)\mapsto(-y,-x)$ 得 $(-y)(-x)=1$；所以 $y=\pm x$ 都是对称轴。取曲线上点 $(1,1)$，关于 $x=0$ 对称为 $(-1,1)$，乘积为 $-1\ne1$，故 $x=0$ 不是对称轴。$x=0,y=0$ 只是两条渐近线。

\textbf{结论：} \textbf{两条渐近线的角平分线是本题的对称轴，但渐近线本身不是。}

\medskip
\textbf{\textcolor{CExample}{例 3（抛物线：一行判定，定义域要对应）}}

\textbf{题目：} 判断 $y=(x-2)^2+1$ 是否关于 $x=2$ 对称。

\textbf{解：} 设 $f(x)=(x-2)^2+1$，则
\[
  f(2+t)=t^2+1=f(2-t)\qquad(t\in\mathbb R).
\]
左右两侧横坐标关于 $2$ 对称，高度相等，故 $\boxed{x=2}$ 是对称轴。

\textbf{结论：} 函数图像关于 $x=h$ 对称，优先检查 $f(h+t)=f(h-t)$。

\medskip
\textbf{\textcolor{CExample}{例 4（平移反比例函数：无需复杂反射计算）}}

\textbf{题目：} 求曲线 $(x-2)(y-3)=1$ 的两条对称轴，并用原方程验证。

\textbf{解：} 设 $u=x-2,v=y-3$，曲线变成 $uv=1$，故候选对称轴是 $v=u$ 与 $v=-u$，分别得到
\[
  \boxed{y=x+1},\qquad\boxed{y=-x+5}.
\]
对第一条轴 $y=x+1$，$P(x,y)$ 的对称点是 $P'(y-1,x+1)$，于是
\[
  (x'-2)(y'-3)=(y-3)(x-2)=1.
\]
对第二条轴 $y=-x+5$，对称点是 $P''(5-y,5-x)$，于是
\[
  (x''-2)(y''-3)=(3-y)(2-x)=1.
\]
所以两条候选轴都成立。

\textbf{结论：} \textbf{平移看结构可以寻找轴；代换原方程才能严谨确认。}

\medskip
\textbf{\textcolor{CExample}{例 5（标准双曲线：渐近线角平分线）}}

\textbf{题目：} 判断 $\dfrac{x^2}{9}-\dfrac{y^2}{4}=1$ 的对称轴与渐近线的关系。

\textbf{解：} 曲线的渐近线为 $y=\frac23x$、$y=-\frac23x$。对称代换 $x\mapsto-x$ 或 $y\mapsto-y$ 都不会改变 $\dfrac{x^2}{9}-\dfrac{y^2}{4}$，因此 $\boxed{x=0}$、$\boxed{y=0}$ 都是对称轴；两条坐标轴正好是两渐近线的角平分线。\textbf{不能把渐近线也当成对称轴。} 例如 $(3,0)$ 在曲线上，关于 $y=\frac23x$ 对称的点为 $(15/13,36/13)$，其横坐标绝对值小于 $3$，不可能位于这条双曲线上。

\textbf{结论：} \textbf{真正的双曲线}可从渐近线角平分线快速\emph{找到}两条轴，再由标准方程的平方项\emph{证明}。

\medskip
\textbf{\textcolor{CExample}{例 6（任意斜直线：代入通式）}}

\textbf{题目：} 不借助平移，验证 $(x-2)(y-3)=1$ 关于 $x-y+1=0$ 对称。

\textbf{解：} 取 $A=1,B=-1,C=1$，有
\[
 s=\frac{x-y+1}{2},\quad
 x'=x-2s=y-1,\quad y'=y+2s=x+1.
\]
代入曲线方程，
\[
  (x'-2)(y'-3)=(y-3)(x-2)=1.
\]
对任意曲线点都成立，又因反射两次返回原点，故所给直线确为对称轴。

\textbf{结论：} 不认识图像也能用 $Ax+By+C=0$ 的通式完成严格验证。

\medskip
\textbf{\textcolor{CExample}{例 7（最小反例：两条渐近线不保证轴对称）}}

\textbf{题目：} 曲线 $y=\dfrac1x+\dfrac1{x^2}$（$x\ne0$）有渐近线 $x=0$、$y=0$。它必关于 $y=x$ 或 $y=-x$ 对称吗？

\textbf{解：} 当 $x$ 趋近 $0$，曲线沿竖直方向无限延伸；当 $|x|$ 越来越大，$y$ 越来越接近 $0$，所以确实有这两条渐近线。取 $P(1,2)$，因为 $1+1=2$，$P$ 在曲线上。关于 $y=x$ 对称得 $(2,1)$，但 $f(2)=\frac12+\frac14=\frac34\ne1$，故 $y=x$ 不是对称轴。关于 $y=-x$ 对称得 $(-2,-1)$，但 $f(-2)=-\frac12+\frac14=-\frac14\ne-1$，故 $y=-x$ 也不是对称轴。

\textbf{结论：} \textbf{渐近线可以提供候选轴，不能代替整条曲线的验证；一个反例点就能排除候选轴。}
\end{examplebox}
